\section{Background}
\label{sec:background}

\textbf{Dec-POMDPs} We consider decentralized partially observable Markov Decision Processes (Dec-POMDPs) \cite{nair2003taming} $G$ with state $s$, action $a$, observation function $\Omega^i(s)$ for each player $i$ and transition function $\mathcal{T}(s, a)$. Players receive a common reward $R(s, a)$. We denote the historical trajectory as $\tau =(s_1, a_1, ..., a_{t-1}, s_t)$ and \textit{action-observation history} (AOH) for player $i$ as $\tau^i = (\Omega^i(s_1), a_1, ..., a_{t-1}, \Omega^i(s_t))$, which encodes the trajectory from player $i$'s point of view. A policy $\pi^i(a|\tau_t^i)$ for player $i$ takes as input an AOH and outputs a distribution over actions. We denote the joint policy as $\pi$. We remark that while the value $V^\pi(\tau)$ of policy $\pi$ from any state $s$ is well-defined, the value of playing $\pi$ from an AOH $\tau^i$ is not well-defined without specifying what policy was played \textit{leading} to $\tau^i$, because it affects the distribution of world states at $\tau^i$. The distribution of world states is often referred to as a belief $\mathcal{B}_{\pi}(\tau|\tau^i)=P(\tau|\tau^i,\pi)$. 

In this work, we restrict ourselves to \textit{bounded-length, turn-based Dec-POMDPs}. Formally, we assume that $G$ reaches a terminal state after at most $t_{max}$ steps, and that only a single agent acts at each state $s$ (other agents have a null action). See Section \ref{sec:turn_based} for a discussion of turn-based vs. simultaneous-action environments. 

\textbf{Deep MARL} has been successfully applied in many Dec-POMDP settings~\citep{mackrl,foerster2019bayesian}. The specific OBL algorithm introduced later uses Recurrent Replay Distributed Deep Q-Networks (R2D2)~\citep{r2d2} as its backbone. In deep Q-learning~\citep{dqn} the agent learns to predict the expected total return for each action given the AOH, $Q(\tau^{i}_{t}, a) = E_{\tau \sim P(\tau | \tau^{i})} R_{t}(\tau)$. $R_{t}(\tau) = \sum_{t'=t}^{\infty} \gamma^{(t' - t)} r_{t'}$ is the forward looking return from time $t$ where $r_{t'}$ is reward at step $t'$ with $a_t = a$ and $a_{t'} = \argmax_{a'} Q(\tau_{t'}, a')$ for $t' > t$ and $\gamma$ is an optional discount factor. The state-of-the-art R2D2 algorithm incorporates many modern best practices such as double-DQN~\citep{double-dqn}, dueling network architecture~\citep{dueling-dqn}, prioritized experience replay~\citep{prioritized-replay}, distributed training setup with parallel running environments~\citep{apex} and recurrent neural network for handling partial observability. 

The straightforward way to apply deep Q-learning to Dec-POMDP settings is Independent Q-Learning (IQL)~\citep{tan93multi} where each agent treats other agents as part of the environment and learns an independent estimate of the expected return without taking other agents' actions into account in the bootstrap process. Many methods~\citep{vdn, qmix} have been proposed to learn joint Q-functions to take advantage of the centralized training and decentralized control structure in Dec-POMDPs. In this work, however, we use the IQL setup with shared neural network weights $\bm{\theta}$ for simplicity and note that any potential improvements to the RL algorithm used are orthogonal to our contribution. 


\textbf{Zero-Shot Coordination.} The most common problem setting for learning in Dec-POMDPs is \emph{self-play} (SP) where a team of agents is trained and tested together. Optimal SP policies typically rely on \emph{arbitrary conventions}, which the entire team can jointly coordinate on during training. However, many real-world problems require agents to coordinate with other, unknown AI agents and humans at test time.
This desiderata was formalized as the \emph{Zero-Shot Coordination} (ZSC) setting by \citet{hu2020other}, where the goal is stated as finding algorithms that allow agents to coordinate with \emph{independently trained} agents at test time, a proxy for the the independent reasoning process in humans. ZSC immediately rules out \emph{arbitrary conventions} as optimal solutions and instead requires learning algorithms that produce robust and, ideally, unique solutions across multiple independent runs. 
